(1) 置\(\chi(t)=A\sin(\omega t)/(\hbar\omega)\)，取\(U_0(t)=e^{-i\chi(t)\sigma_z}\)。将原态写\(|\psi\rangle=U_0|\widetilde\psi\rangle\)，变换哈密顿量为\(\widetilde H=U_0^\dagger HU_0-i\hbar U_0^\dagger\dot U_0\)。因\(-i\hbar U_0^\dagger\dot U_0=-A\cos(\omega t)\sigma_z\)，它与原纵向驱动严格相消。
(2) 使用\(\sigma_\pm=(\sigma_x\pm i\sigma_y)/2\)与\(U_0^\dagger\sigma_\pm U_0=e^{\pm 2i\chi(t)}\sigma_\pm\)，得\(\widetilde H(t)=\frac\Delta2[e^{2i\chi(t)}\sigma_++e^{-2i\chi(t)}\sigma_-]\)。Jacobi–Anger公式\(e^{iz\sin\omega t}=\sum_{n=-\infty}^\infty J_n(z)e^{in\omega t}\)将耦合分解成多光子边带，驱动强度参数是\(z=2A/(\hbar\omega)\)。
(3) 若\(\Delta/(\hbar\omega)\ll1\)，保留Magnus/Floquet展开首个周期平均项。因为\(\langle e^{\pm iz\sin\omega t}\rangle=J_0(z)\)，故\(H_{\rm eff}^{(0)}=\frac\Delta2J_0(2A/(\hbar\omega))\sigma_x\)。准能级的首阶有效劈裂为\(|\Delta_{\rm eff}|=|\Delta J_0(2A/(\hbar\omega))|\)，不同于简单把静态\(\Delta\)替换为\(\Delta-2A\)。
(4) 当\(2A/(\hbar\omega)\)取\(J_0\)的零点，例如首根\(2.404826\ldots\)，首阶有效隧穿被抑制，称为相干隧穿抑制。该零点不是有限频率下严格零耦合的普遍定理：高阶Magnus项、非正弦驱动、噪声和与其他能级的多光子共振会造成残余劈裂，必须核查\(\hbar\omega\gg\Delta\)及二能级近似。
